\section{Related Work}

\paragraph{Masked autoencoding for LiDAR.}
The MAE paradigm~\cite{mae} was extended to LiDAR by
Occupancy-MAE~\cite{occmae} and Voxel/LiDAR-MAE~\cite{hess2023lidar}, with
MAELi~\cite{maeli} scaling it to large sequences. Occupancy-MAE introduced
range-aware masking and reconstructs binary voxel occupancy. The proposed
OMU-MAE keeps its range-aware schedule but \emph{adds} a second target at
masked positions---a $768$-d vector matching the frozen DINOv2 patch
feature projected from the paired image. The joint dual objective
(occupancy plus frozen-feature reconstruction, both restricted to masked
positions) is proposed in this work; the occupancy term inherits
Occupancy-MAE's geometry supervision while the feature term injects
semantics that a binary target cannot carry. Regressing \emph{frozen}
teacher features at masked positions is well validated in 2D---MaskFeat's
ablations~\cite{maskfeat}, MILAN~\cite{milan} and EVA~\cite{eva} all show
it outperforms raw-pixel targets; OMU-MAE lifts this recipe to
cross-modal voxel-level driving scenes.

\paragraph{Multi-modal masked autoencoders.}
UniM\textsuperscript{2}AE~\cite{unim2ae} jointly trains both modality
encoders and reconstructs \emph{raw} pixels and LiDAR through a 3D
interaction module; NS-MAE~\cite{nsmae} reconstructs raw observations via
NeRF-style volumetric rendering. OMU-MAE differs on three axes, each with
a concrete payoff: the image encoder is \emph{frozen}, which removes the
collapse risk of joint training and halves the trainable footprint; the
cross-modal target is the frozen DINOv2 \emph{feature} volume rather than
raw pixels, so supervision carries semantics instead of photometric
nuisance; and there is no interaction module or back-projection---only
direct concatenation into the voxel input---which keeps the whole pipeline
a single 3D CNN branch that trains in a few hours on one GPU. The
trade-off is intentional: less expressive than a jointly trained
two-branch model, but far cheaper and collapse-free by construction.

\paragraph{Vision-foundation-model distillation to LiDAR.}
SLidR~\cite{slidr}, ScaLR~\cite{scalr} and SuperFlow++~\cite{superflow}
distill frozen-image features into LiDAR superpoints via InfoNCE at sparse
superpoint resolution; CleverDistiller~\cite{cleverdistiller} replaces
InfoNCE with a direct feature-similarity loss plus an occupancy auxiliary.
The frozen-VFM target $+$ occupancy task is shared with OMU-MAE; the
differentiator is \emph{masked reconstruction}. To make this comparison
concrete we re-implement both a contrastive (SLidR) and a no-mask
distillation (CleverDistiller-equivalent) variant inside our framework and
compare under an identical protocol (Sec.~\ref{sec:setup}).

\paragraph{JEPA and dense vision features.}
JEPA-family methods~\cite{ijepa,vjepa,lejepa,adljepa} predict latent
representations of masked targets, but their co-evolving target requires
careful regularization (VICReg~\cite{vicreg}) to avoid collapse; we
deliberately sidestep this with a frozen target. We use
DINOv2~\cite{dino,dinov2,dinov3} ViT-B/14 as that target for the quality of
its dense features and the public availability of weights via
\texttt{torch.hub}.
